\section{A global game of stablecoin redemption}\label{sec:model}

This section formalises the depeg channel of Section~\ref{sec:risk_factors} as a global game and states the equilibrium properties the calibration relies on. We use a one-class representative-holder specification; the multi-class generalisation collapses to it when the population concentrates on a single class, and a two-class robustness extension is reported in the internet appendix. Proofs and the numerical solver are in Appendix~\ref{app:proofs}.

\subsection{Primitives and information}\label{sub:model_primitives}

There is a continuum of homogeneous holders of unit mass. The latent issuer fundamental is $\theta \in \mathbb{R}$, drawn from a normal prior $\mathcal{N}(\mu_\theta, 1/\tau)$ with regime-dependent mean and precision; the regime $\rho \in \{b, s\}$ is drawn from a Bernoulli with $\mathbb{P}(\rho=s) = p_s$. In the baseline regime $\mu_\theta = \mu_b$ and $\tau = \tau_b$, fixed at $\mu_b = 1$ and $\tau_b = 400$ so the baseline is quasi-deterministic at par; the stress-regime $\mu_s, \tau_s$ are estimated. Each holder $i$ observes a private signal
\begin{equation}\label{eq:signal}
s_i = \theta + \sigma \varepsilon_i, \qquad \varepsilon_i \sim \mathcal{N}(0,1),
\end{equation}
with $\sigma > 0$ the signal noise, so the posterior is $\theta \mid s_i \sim \mathcal{N}(s_i, \sigma^2)$ \citep{morris_shin_1998}. The issuer holds primary reserves $R$ and asset value $A$, with default hazard $h \in [0,1]$ and loss given default $\xi \in [0,1]$. The secondary market has baseline depth $M$, price-impact scale $\psi > 0$, convexity $\nu > 1$, congestion penalty $\mu_q$, and depth-erosion $\zeta$. The per-redemption transaction cost is $\phi$ and the convenience yield of holding the coin is $\omega$.

\subsection{Redemption and the secondary price}\label{sub:model_secondary}

Each holder chooses to redeem or hold. By symmetry and the monotone-likelihood-ratio property of the signal, the equilibrium is in threshold strategies, redeem if and only if $s_i < s^*$, so the mass that redeems at fundamental $\theta$ is
\begin{equation}\label{eq:D}
D(\theta; s^*) = \Phi\!\left(\frac{s^* - \theta}{\sigma}\right),
\end{equation}
with $\Phi$ the standard normal CDF. Redemption runs through two channels. Each holder first attempts primary redemption at the issuer, paid at par up to liquid reserves $R$; the unmet demand spills into the secondary market,
\begin{equation}\label{eq:U}
U(D) = \max\{D-R, 0\},
\end{equation}
which clears at
\begin{equation}\label{eq:psec}
p^{sec}(D) = \min\!\left\{1, \, \max\!\left[\underline{p}, \, 1 - \psi\!\left(\frac{U(D)}{\widetilde M(D)}\right)^\nu - \mu_q q\right]\right\}, \qquad
\widetilde M(D) = \frac{M}{1 + \zeta\, U(D)/M},
\end{equation}
where $\widetilde M(D)$ is the depth that itself erodes as unmet demand grows, $\underline p$ a price floor, and $q$ the rail queue ratio (zero in the specification used here). The price falls convexly ($\nu>1$) in unmet redemption relative to depth, the feature the mechanism check of Section~\ref{sub:mechanism_check} tests.

\subsection{Payoffs and equilibrium}\label{sub:model_equilibrium}

Conditional on $\theta$ and the aggregate $D$, the payoff to redeeming is
\begin{equation}\label{eq:VR}
V_R(D) = \min\!\left\{1, \tfrac{R}{D}\right\} + \left(1 - \min\!\left\{1, \tfrac{R}{D}\right\}\right) p^{sec}(D) - \phi,
\end{equation}
par on the share met from reserves and the depressed secondary price on the rest, and the payoff to holding is
\begin{equation}\label{eq:VH}
V_H(\theta) = \theta - c_q q + \omega - \xi h.
\end{equation}
The issuer counterparty loss $\xi h$ enters the value of \emph{holding}, not the realised price. This placement is the substantive modelling choice: a counterparty loss netted from the realised price would be common to redeeming and holding and would cancel in the indifference condition below, leaving the threshold unmoved. In $V_H$ it lowers the value of holding, raises the conversion mass, and feeds the coordination channel, so a heavier counterparty profile both adds a loss component and tightens the run incentive; at the calibrated parameters this second, coupling effect is second-order relative to the additive component (Section~\ref{sec:identification}). The marginal holder at $s_i = s^*$ is indifferent in expectation over $\theta$ given the posterior $\mathcal{N}(s^*, \sigma^2)$:
\begin{equation}\label{eq:F}
F(s^*) := \mathbb{E}_{\theta \mid s_i = s^*}\!\left[V_R(D(\theta; s^*)) - V_H(\theta)\right] = 0.
\end{equation}
This is the endogenous fixed point. The threshold $s^*$ enters $F$ both as the conditioning event of the posterior and through the aggregate mass $D(\theta; s^*)$ that all holders applying the same threshold imply; the two-way dependence is the run, and it admits no closed form.

\subsection{Existence, uniqueness, and comparative statics}\label{sub:model_propositions}

Three propositions give the model the properties the calibration relies on; all are proved in Appendix~\ref{app:proofs}.

\begin{proposition}[Existence and uniqueness]\label{prop:existence}
Assume $\nu > 1$, $\sigma > 0$, $0 < R < 1$, $M>0$, and that the best-response operator $T$ on threshold profiles is a contraction in the sup norm. Then \eqref{eq:F} has a unique solution $s^* \in \mathbb{R}$.
\end{proposition}

The proposition gives a single equilibrium threshold $s^*$ at any admissible parameter vector, so the conversion mass $D^*(\theta)$ is determined at each fundamental and the calibration recovers one expected loss per issuer. The contraction condition is verified numerically at the calibrated parameters.

\begin{proposition}[Reserve comparative static]\label{prop:reserves}
At the equilibrium $s^*$, $\partial D^*/\partial R \geq 0$ in the strict interior of the parameter space.
\end{proposition}

A deeper liquid reserve sustains a higher secondary price at any conversion volume, which raises the value of converting and so weakly raises the equilibrium conversion mass. This is the reserve-quality channel of \citet{ma_zeng_zhang_2025}. The sign reads as the opposite of the bank-run intuition, and the difference is worth stating rather than leaving implicit. In \citet{goldstein_pauzner_2005_demand_deposits} stronger fundamentals shrink the run region because withdrawal is destructive: the deposit contract liquidates assets at a loss and late claimants stand behind early ones, so safety means fewer agents running. Redemption here is not destructive withdrawal but conversion at a price. More par capacity makes exit safer, so weakly more holders take it, and what deeper reserves reduce is not the volume of conversion but its cost: a larger share clears at par and the spillover clears against a higher secondary price. The conversion mass $D^*$ is therefore not the policy object; the priced loss is. Numerically, at the calibrated parameters the expected loss falls monotonically in $R$ while the conversion mass rises: re-solving the model on a reserve grid from $R=0.4$ to $R=0.8$ moves $\Pi_{\text{risk}}$ from 63.4 to 47.0 basis points for USDC, with the conditional stress depeg falling from 462 to 102 basis points, and from 71.3 to 71.2 for USDT; USDC's baseline conversion mass rises from 0.485 to 0.497 over the same grid. So $\partial \Pi_{\text{risk}}/\partial R < 0$ holds throughout the grid even as $\partial D^*/\partial R \geq 0$, and a supervisor should read Proposition~\ref{prop:reserves} jointly with that comparative static: liquidity requirements lower the holder's expected loss precisely by making redemption safe enough that more of it happens at par.

\begin{proposition}[Convenience-yield comparative static]\label{prop:convenience}
At the equilibrium $s^*$, $\partial D^*/\partial \omega < 0$ in the strict interior of the parameter space.
\end{proposition}

A higher payment-use convenience yield raises the value of holding and so lowers the equilibrium conversion mass: a more useful stablecoin retains holders even at marginally lower fundamentals. This is the adoption-fragility channel of \citet{bertsch_2025_adoption_fragility}.

Together the two comparative statics formalise the claim of Section~\ref{sub:not_homogeneous} that reserve composition and payment use are primitives that move the priced expected loss, not stylistic differences between issuers.

\subsection{The expected loss}\label{sub:what_it_delivers}

The output of the model, for each issuer under its own calibration, is three expected per-unit losses in basis points,
\begin{align}
\ell_C &= h \cdot T \cdot \xi \cdot 10^4, \qquad T = 1,\label{eq:ellC} \\
\ell_{D,b} &= \mathbb{E}_b\!\left[D(\theta; s^*)\,(1 - p^{sec}(D(\theta; s^*)))\right] \cdot 10^4,\label{eq:ellDb} \\
\ell_{D,s} &= \mathbb{E}_s\!\left[D(\theta; s^*)\,(1 - p^{sec}(D(\theta; s^*)))\right] \cdot 10^4,\label{eq:ellDs}
\end{align}
where $\mathbb{E}_\rho$ is the expectation under regime $\rho$. The counterparty component $\ell_C$ is the reduced-form credit term of Section~\ref{sub:counterparty}, evaluated at the one-year exposure window $T=1$ of the continuous-balance holder defined there and independent of the coordination stage; the two depeg components are the structural output of the global game. Their probability-weighted sum is the expected loss the holder bears,
\begin{equation}\label{eq:total_risk_premium}
\Pi_{\text{risk}} = \ell_C + (1-p_s)\,\ell_{D,b} + p_s\,\ell_{D,s},
\end{equation}
which we add to the on-chain transfer fee. Section~\ref{sec:calibration} calibrates this expected loss separately for the two stablecoins.
